\documentclass[11pt]{article}
\usepackage[a4paper,margin=1in]{geometry}
\usepackage{amsmath,amssymb,amsthm,booktabs,array}
\usepackage[hidelinks]{hyperref}
\makeatletter
\g@addto@macro\UrlBreaks{\do\-}
\makeatother
\pdftrailerid{}

\newtheorem{theorem}{Theorem}
\newtheorem{proposition}[theorem]{Proposition}
\newcommand{\Q}{\mathbb Q}
\newcommand{\Cl}{\operatorname{Cl}}
\newcommand{\Norm}{\operatorname{N}}

\title{A Certified Canonical Census of One-Place Stark Invariants over Real Quadratic Fields\\
\large Exact Quadratic Value Orbits and the Higher-Order Frontier}
\author{Hainan Zhao\\\small Independent Researcher}
\date{Version 1.1 correction candidate, 1 August 2026\\
\small Version 1.1: \url{https://doi.org/10.5281/zenodo.21730707}}

\begin{document}
\maketitle

\begin{abstract}
We give a certified finite census of canonically selected one-place
differenced Stark packets over real quadratic fields.  For squarefree
radicands $2\leq D\leq200$ and finite-ideal norm at most $100$, we first
quotient finite ideals by field conjugation, choose the least normalized
HNF in each orbit, and then attach the pinned real place $\infty_2$.
This gives $8{,}200$ selected moduli, which split exactly into $3{,}936$
empty-support rows, $1{,}560$ quadratic-support rows, and $2{,}704$
higher-order rows.  Every quadratic row has an exact squarefree
polynomial for its distinct packet-value orbit over the real quadratic
base.  A deleted-prime cover criterion decides, before any regulator or
unit computation, exactly when the whole quadratic packet is one.  The
higher-order table has complete registered-mechanism status
for $2{,}699$ rows and preserves five tool-incomplete quartic rows.
These are statements about the selected finite universe only: we claim
neither a census of all pointed one-place moduli, an asymptotic density,
nor a new higher-order packet identity.
\end{abstract}

\section{Claim boundary and conventions}

Let $K$ be real quadratic and $\mathfrak m=\mathfrak f\infty_2$ a
one-place modulus.  With $R$ the real-place sign class, put
\[
 Z_{\mathfrak m}(s,A)=\zeta_{\mathfrak m}(s,A)-
 \zeta_{\mathfrak m}(s,RA),\qquad
 X_A=\exp Z'_{\mathfrak m}(0,A).
\]
Here $\Cl_{\mathfrak m}(K)$ is the ray class group for the indicated
finite and infinite modulus.  The packet is the Artin-indexed tuple
$(X_A)_{A\in\Cl_{\mathfrak m}(K)}$.
The exact Fourier, Artin-action, place-ordering, and positive-orientation
conventions are those of the companion results paper \cite{ZhaoResults};
this paper does not introduce a second convention system.

The proved assertions here are the finite trichotomy, the exhaustive
quadratic packet-value-orbit corpus, the quadratic deleted-prime cover
criterion, and twenty-two named higher-order member transports.  The
higher-order eligibility table is an exact
finite-population observation: five quartic rows remain incomplete.
We do not assert results beyond the frozen selected range.  Eligibility
is not a packet identity.  In particular, a common normal closure or a
banked representative does not promote another member: exact conductor,
ray-class, orientation, and Artin-labelled packet transport are separate
obligations.  The registered Shintani-transfer scope has 232 eligible
rows.  Its successor ledger records twenty-two completed noncanonical member
transports; the other 210 remain open.

We use \emph{PROVED} when a cited theorem applies after exact checks and
\emph{OBSERVED} for reproducible descriptive statistics of this finite
population.  The tag \texttt{CERTIFIED\_NUMERICAL} denotes a rigorous
enclosure with a printed radius criterion.
``Unconditional'' means that an asserted identity follows from a proved
theorem after exact hypothesis checks; it does not prove the general
Stark conjecture.

\section{Frozen universe and trichotomy}

The selected universe consists of maximal-order fields $\Q(\sqrt D)$
for squarefree $2\leq D\leq200$ and nonzero integral ideals
$\mathfrak f$ with $\Norm\mathfrak f\leq100$.  In the frozen integral
basis, replace $\mathfrak f$ by the lexicographically smaller normalized
HNF of $\mathfrak f$ and $\overline{\mathfrak f}$, then attach the
pinned place $\infty_2$.  Thus conjugation is applied to the finite
ideal before the distinguished place is fixed.  This is a canonical
selected-modulus census; it is not the full isomorphism-class quotient
of all pairs $(\mathfrak f,\infty_i)$.

A clean PARI replay enumerates 121 fields and 13,939 raw ideals.  Of
these, 2,461 are self-conjugate and 11,478 are nonself-conjugate, so the
implemented finite-ideal quotient has
\begin{equation}\label{eq:universe-count}
 2461+5739=8200
\end{equation}
representatives.  No theorem below promotes this selected universe to
a census of all pointed one-place moduli.

Let $T$ be empty differenced Fourier support; let $Q$ be nonempty
support in which every supported character has order two; and let $H$
be nonempty support containing a character of order greater than two.
The structural sign-class lemmas in \cite{ZhaoResults} show that empty
support implies $X_A=1$ for every $A$.  The converse is false: a
nonempty supported term can be annihilated by an imprimitive Euler
factor.

\begin{theorem}[Frozen trichotomy]\label{thm:trichotomy}
In the frozen universe,
\[
 |T|=3936,\qquad |Q|=1560,\qquad |H|=2704.
\]
Every row belongs to exactly one stratum.  Every packet in $T$ is one.
\end{theorem}

\begin{proof}
The clean enumerator certifies every quadratic base field and applies
the canonicalization in \eqref{eq:universe-count}.  Exact character
inversion computes the differenced support of every retained row.
The three predicates above are disjoint and exhaustive; their audited
counts are $3936$, $1560$, and $2704$.  If the support is empty, Fourier
inversion gives $Z_{\mathfrak m}(s,A)=0$ identically, hence $X_A=1$.
\end{proof}

The earlier v5 routing declaration recorded 3,899 trivial rows.  The
remaining 37 are empty-support rows that its obsolete routing order
assigned an exponent-cap label before applying the empty-support theorem.
The historical record is preserved; Theorem~\ref{thm:trichotomy} uses the
support-first replay.

The set $T$ is therefore a structural empty-support stratum, not the
set of all rows whose evaluated packet is one.  In particular, the
quadratic stratum below contains a disjoint 346-row value-one
degeneracy sub-stratum.

\section{The exact quadratic stratum}

For a supported quadratic character, the uniform theorem in
\cite{ZhaoResults}, based on Tate's rank-one algebraicity theorem
\cite[Thm.~IV.5.4]{Tate1984}, gives exact exponents from imprimitive Euler
factors and a relative regulator index.  Conductors, Euler factors,
relative-unit kernels, orientations, and indices are checked exactly.
There are 2,232 supported character occurrences: 672 have zero Euler
factor, and 346 rows have all factors zero and packet polynomial $X-1$.

More explicitly, for $G=\Cl_{\mathfrak m}(K)$ the theorem gives
\begin{equation}\label{eq:quadratic-product}
 X_A=\prod_{\chi(R)=-1}|\tau_1u_\chi|^{c_{A,\chi}},\qquad
 c_{A,\chi}=\frac{4}{|G|}\chi(A)^{-1}E_\chi
 \frac{h_{L_\chi}}{h_K}\frac{w_K}{w_{L_\chi}}\frac1{I_\chi}.
\end{equation}
Every quantity in the exponent and every oriented relative unit is
computed by exact arithmetic.

If $u_i$ is a norm-one relative unit, then
\[t_i=u_i+u_i^{-1}=\operatorname{Tr}_{L_i/K}(u_i)\in K.\]
The formal sign-orbit polynomial is therefore synthesized without its
compositum by
\begin{equation}\label{eq:trace-descent}
P_0(x)=x-1,\qquad
P_i(x)=\operatorname{Res}_z\bigl(P_{i-1}(z),x^2-t_i xz+z^2\bigr).
\end{equation}
Each root $z$ is replaced by $zu_i$ and $zu_i^{-1}$, so every intermediate
polynomial remains over $K$.  The recurrence concerns
denominator-cleared packet powers.  At denominator two, the positive
packet factor is selected only after exact square-lift, reciprocity,
irreducibility, positivity, and Artin-image checks.

For a quadratic row define its \emph{packet-value-orbit polynomial} by
\[
 F_{\mathfrak m}(x)=\prod_{v\in\{X_A:A\in G\}_{\rm distinct}}(x-v).
\]
This squarefree polynomial records the distinct effective Artin orbit;
it is not, by itself, a class-to-root labelling.  The exact ray-group
coordinates, sign-class log, supported characters, and exponents in
\eqref{eq:quadratic-product} are retained so that the Artin-labelled
tuple can be reconstructed on replay.  Only the worked row below prints
that full labelling in the present article.

\begin{theorem}[Exhaustive quadratic stratum]\label{thm:quadratic}
Every one of the 1,560 rows of $Q$ has a certified exact squarefree
packet-value-orbit polynomial $F_{\mathfrak m}(x)$ over its base field.
Its degree equals the exact effective Artin-image size.  The stable
RQ-ordered row certificates have terminal hash-chain root
\begin{center}\small\texttt{7c04242b1d4c11293af96f83f4915dbed25f6125c60d82965e533df5c9d81855}.\end{center}
\end{theorem}

\begin{proof}
Apply \eqref{eq:quadratic-product} character by character.  Clear the
common rational denominator, synthesize the formal sign orbit by
\eqref{eq:trace-descent}, and restrict it to the exact Artin sign image.
For the 69 denominator-two rows, exact squareness in the full ray field
selects the positive lift.  Each retained factor is then checked for
reciprocity, squarefreeness, irreducibility over $K$, positive real-root
signature, and equality of its degree with the effective Artin-image
size.  All 1,560 row checks pass.  The hash chain protects the ordered
corpus against later alteration; it is an integrity record, not a
substitute for these algebraic checks.
\end{proof}

For a supported quadratic character $\chi$, let $\chi^\circ$ be its
primitive character, let $\mathfrak f_\chi$ be its finite conductor,
and put
\[
 \mathcal D_\chi=\{\mathfrak p:\mathfrak p\mid\mathfrak f,
                    \ \mathfrak p\nmid\mathfrak f_\chi\}.
\]
These are the prime ideals deleted when the primitive $L$-function is
made imprimitive at $\mathfrak f$.

\begin{theorem}[Deleted-prime cover criterion]\label{thm:cover}
In the quadratic-support regime, the following are equivalent:
\begin{enumerate}
\renewcommand{\theenumi}{\roman{enumi}}
\renewcommand{\labelenumi}{(\theenumi)}
\item $X_A=1$ for every $A\in G$;
\item $L'_{\mathfrak m}(0,\chi)=0$ for every $\chi(R)=-1$;
\item $E_\chi=0$ for every $\chi(R)=-1$;
\item for every $\chi(R)=-1$ there is a
      $\mathfrak p\in\mathcal D_\chi$ with
      $\chi^\circ(\mathfrak p)=1$.
\end{enumerate}
Equivalently, the deleted primes that split in the quadratic extension
cut out by $\chi^\circ$ cover every supported character.
\end{theorem}

\begin{proof}
The imprimitive factor is
\[
 E_\chi=\prod_{\mathfrak p\in\mathcal D_\chi}
              (1-\chi^\circ(\mathfrak p)).
\]
Every displayed local value is $1$ or $-1$, because $\chi^\circ$ is
quadratic and unramified at the deleted prime.  Hence $E_\chi=0$ exactly
when one such value is $1$, proving the equivalence of (iii) and (iv).
The companion paper proves the explicit formula and its regulator
identity in \cite[Prop.~4 and Eq.~(9)]{ZhaoResults}.  More explicitly,
$u_\chi$ generates the primitive rank-one norm kernel modulo torsion,
so it is non-torsion; injectivity of the logarithmic embedding on the
free unit group makes its relative vector $(b,-b,0)$ have $b\ne0$, and
the pinned orientation makes $b=\log|\tau_1u_\chi|>0$.  The class
numbers, roots-of-unity quotient, and index are also nonzero.  Thus the
explicit quadratic formula gives (ii) if and only if (iii).  Fourier
inversion gives (ii) implies (i).  Conversely, (i) gives
$Z'_{\mathfrak m}(0,A)=0$ for every $A$; character orthogonality applied
to the exact inverse transform recovers (ii).
\end{proof}

The criterion is independent of the census bounds and is an exact
pre-regulator decision test.  In the frozen range its deleted-prime
export checks 2,232 supported character occurrences and 1,516 deleted
prime occurrences.  Exactly 699 deleted occurrences have local value
$1$, and the resulting cover holds on exactly 346 rows: 296 with one
supported character and 50 with two.  It has zero character-level and
row-level classification errors, and agrees with $F_{\mathfrak m}=x-1$
on all 1,560 exact packet polynomials.

No four-support row in the frozen range is covered, but this is not a
general theorem.  A preregistered falsification search finds the exact
counterexample $K=\Q(\sqrt7)$ and finite-ideal HNF $[42,0;0,6]$ of norm
252.  Its ray group is $C_6\times C_2\times C_2$, its sign log is
$(3,1,0)$, and its four supported characters all have a deleted prime
with local value $1$.  Thus all four Euler products vanish and the
packet is one.  This row lies outside the frozen census and refutes the
tempting universal ``four-support nondegeneracy'' pattern.

\begin{table}[ht]
\centering
\caption{Exact quadratic-stratum distributions.}\label{tab:quadratic}
\begin{tabular}{lr}\toprule
invariant & count\\\midrule
packet degree $1,2,4,8$ over $K$ & $346,930,242,42$\\
common denominator $1,2$ & $1491,69$\\
all-Euler-zero packet $X-1$ & $346$\\
maximum coefficient-coordinate size & $62$ digits\\\bottomrule
\end{tabular}
\end{table}

Here the coefficient-coordinate size is the maximum decimal digit count
of the numerator or denominator of either coordinate in the frozen
basis $1,y$ of $K$.  The preregistered height-only predictor had maximum
89 digits, below the frozen 256-digit cap.  The full exact run took
43.46 seconds under PARI/GP 2.15.4; this is an OBSERVED engineering
statistic, not a portable benchmark.

A separately frozen 50-row audit used the independent analytic
class-number/regulator formula before opening the corpus traces.  At
384 bits all 43 nonzero character comparisons and 101 Artin sign rows
pass the preregistered $10^{-38}$ difference-radius cap; the independent
balls were inflated by $10^{-45}$.  This is a CERTIFIED\_NUMERICAL
consistency check, not a proof input.  The initial 192-bit run failed
the radius gate at RQ-006617; that failure is preserved rather than
discarded.

\subsection{A worked imprimitive row}

RQ-000013 over $\Q(\sqrt2)$ has finite ideal HNF
$[14,6;0,2]$ of norm 28 and ray group $C_2$.  Its supported character
has primitive conductor obtained by removing one prime of norm two.
The exact imprimitive Euler factor and relative index are both two:
\[ E_\chi=2,\qquad I_\chi=2. \]
The oriented relative unit has minimal polynomial
\[u^4-2u^3+u^2-2u+1\]
and its selected real root lies in $(9/5,19/10)$.  Exact Artin labels
then give
\[ X_{[0]}=u^2,\qquad X_{[1]}=u^{-2}. \]
This is a PROVED application of the uniform theorem, including the
otherwise easy-to-miss nonzero imprimitive Euler branch.  A numerical
point-value comparison is retained only as an OBSERVED cross-check.

\section{Higher-order taxonomy and frontier}

For every $H$ row we record the status of three registered mechanisms.
The Shintani-transfer column requires the exact operational predicates
used in the companion theorem: derived-subgroup order two, the
real-place and unit-congruence gates, and the frozen effective-exponent
bound.  The cyclic-quartic CM column requires order-four nonquadratic
support, projective $V_4$ geometry, exact local factors, the registered
$S$-set condition, roots of unity, and orientation.  These two columns
record route eligibility, not packet identities.  They are the
mutually disjoint assignment sets frozen by the v5 declaration: its
exact set-intersection audit finds no doubly assigned row in this
finite corpus.  This is not a general theorem that the two mechanisms
cannot overlap; a future doubly eligible row would require an explicit
overlap category rather than silent priority assignment.

The Roblot column tests Theorem~7.1 \cite{Roblot2013}: conditions
(A1)--(A3), equality with the theorem's $S$-set, class number prime to
three, and absence of wild ramification above three.  Its conclusion is
the theorem's weak cyclic-sextic statement up to absolute values, not an
Artin-labelled positive packet.  The registered-mechanism status is
complete on 2,699 rows; five old quartic constructions remain
tool-incomplete.  Every sextic applicability decision is complete.

\begin{table}[ht]
\centering
\caption{Route-by-Roblot cross-table for the finite $H$ stratum.}
\label{tab:higher}
\begin{tabular}{lrrrrr}\toprule
registered route & weak coverage & hyp. fail & unsupported & incomplete & total\\\midrule
Shintani transfer & 70 & 45 & 117 & 0 & 232\\
cyclic-quartic CM & 782 & 99 & 0 & 0 & 881\\
exclusive frontier & 227 & 261 & 1098 & 5 & 1591\\\midrule
total & 1079 & 405 & 1215 & 5 & 2704\\\bottomrule
\end{tabular}
\end{table}

Thus 1,359 rows---the 261 hypothesis failures and 1,098 unsupported
orders in the exclusive frontier---fail every completed registered
mechanism.  The 227 frontier rows with Roblot weak coverage are not
Artin-labelled packet identifications, and the five incomplete rows are
not counted as failures.

For order-six support, 764 character/inverse-character kernels reduce
to 407 primitive-field keys.  Exact local checks remove 25 keys and
382 require field certificates.  Of those 382, 309 complete the field
gates directly: 48 reuse full \texttt{bnfcertify} certificates and 261
use quotient certificates from computed prime-to-three class groups.
The remaining 73 possible three-divisible class-number fields are each
closed by an exact unramified cyclic-cubic certificate.  Thus every
sextic kernel has an exact Roblot applicability decision.  The five incomplete entries of
Table~\ref{tab:higher} are old quartic construction failures, not
missing sextic decisions and not mathematical no-go results.

The index column uses the genuine derived-subgroup order in the
common stable modulus, rather than the legacy W1 screening proxy.
This distinction resolves a potentially misleading control case:
RQ-005298 has legacy displayed index two, but its certified normal
closure has relative degree $128$, maximal absolutely abelian subfield
of relative degree $32$, and therefore derived-subgroup order
$128/32=4$.  It fails the actual Engine-B index-two predicate; its
non-eligibility is not an unstated failure of a closure, conductor, or
transfer condition.  The separately incomplete quartic Engine-C
resolvent for this row has no mathematical verdict.

The headline wall is RQ-000692 over $\Q(\sqrt{21})$, with finite norm
36 and support profile $(2,6)$.  Its Shintani index is six.  The
one-place, local-$S$, and prime-to-three-class-number conditions of
Roblot's theorem pass, but a prime above 3 is wildly ramified in the
relevant sextic extension.  The theorem therefore does not apply.  This
is an exact limitation of the registered methods, not a non-algebraicity
theorem.

\subsection{Minimal frontier targets}

The frontier artifact selects the least RQ identifier at each support
order, independently of the outcome of later mechanism checks.  A
representative initial segment is given in Table~\ref{tab:frontier}.
The complete table, including orders through 210, is machine-readable
in the H-taxonomy artifact.  The distinction between ``unresolved'' and
``all mechanisms fail'' is retained: an unresolved row can still be
eligible for one of the registered mechanisms.

\begin{table}[ht]
\centering
\caption{First frozen frontier targets by support order.}\label{tab:frontier}
\begin{tabular}{rll}\toprule
order & minimal unresolved & minimal all-mechanisms-fail\\\midrule
4 & RQ-000030 & RQ-000277\\
6 & RQ-000032 & RQ-000227\\
8 & RQ-000018 & RQ-000018\\
10 & RQ-000078 & RQ-000078\\
12 & RQ-000294 & RQ-000294\\
14 & RQ-000524 & RQ-000524\\
16 & RQ-001056 & RQ-001056\\
18 & RQ-000657 & RQ-000657\\
20 & RQ-000785 & RQ-000785\\
24 & RQ-000363 & RQ-000363\\\bottomrule
\end{tabular}
\end{table}

\section{Engine-B transport scope}

The registered Shintani-transfer population consists of 232 rows in 88
groups by exact normal-closure polynomial.  A successor ledger promotes
the first twelve noncanonical members in Table~\ref{tab:transport}; a
corrected integral-basis batch adds ten further transports in the
successor ledger.
For RQ-000039, for example,
\[
 X_{98}(A)=X_{49}(A)X_{49}(A\mathfrak q^{-1})^{-1},
\]
where $\mathfrak q$ is the unique norm-two prime of $\Q(\sqrt2)$.
Every row uses exact Euler deletion and a positive product or quotient
at the frozen split embedding; this is a distribution relation, not
equality of packets.  In the table, $c$ means that target label $a$
maps to source label $ca$ in $C_6$, and $(\ell,e)$ records a deleted
prime with source ray log $\ell$ and quotient exponent $e$.

\begin{table}[ht]
\centering\small
\caption{First twelve proved noncanonical Engine-B member transports.}
\label{tab:transport}
\begin{tabular}{llllll}\toprule
target & source & closure & $c$ & deleted $(\ell,e)$ & status\\\midrule
RQ-000039 & RQ-000021 & B5-015 & 1 & $(1,1)$ & PROVED\\
RQ-000195 & RQ-000190 & B5-025 & 1 & $(1,1)$ & PROVED\\
RQ-000200 & RQ-000190 & B5-025 & 1 & $(2,1)$ & PROVED\\
RQ-000205 & RQ-000190 & B5-025 & 1 & $(1,2)$ & PROVED\\
RQ-000213 & RQ-000190 & B5-025 & 1 & $(1,1),(2,1)$ & PROVED\\
RQ-000221 & RQ-000190 & B5-025 & 5 & $(1,3)$ & PROVED\\
RQ-000228 & RQ-000190 & B5-025 & 5 & $(5,2)$ & PROVED\\
RQ-000425 & RQ-000419 & B5-022 & 1 & $(4,1)$ & PROVED\\
RQ-000436 & RQ-000419 & B5-022 & 5 & $(1,1)$ & PROVED\\
RQ-000457 & RQ-000419 & B5-022 & 5 & $(4,1),(4,1)$ & PROVED\\
RQ-000459 & RQ-000419 & B5-022 & 1 & $(5,1)$ & PROVED\\
RQ-000465 & RQ-000419 & B5-022 & 5 & $(3,1)$ & PROVED\\\bottomrule
\end{tabular}
\end{table}

The earlier B5-086/B5-021/B5-033 direct-source exclusions used
nonintegral quadratic bases with HNFs expressed in the frozen integral
bases, and are withdrawn.  The corrected all-closure screen finds ten
geometrically eligible directions from their certified sources, and
the label-aware Euler and orientation proofs now promote all ten.
Across the former 220 unpromoted members, 116 have no incoming direct
coprime-deletion direction and 104 retained a
source-or-proof-open direction.  These are route classifications, not
negative packet claims.  The remaining 210 members retain status
\texttt{UNSTARTED\_NO\_CASE\_LEVEL\_PACKET\_CLAIM}.

Six historical rows---RQ-000970, RQ-000991, RQ-000993, RQ-001004,
RQ-002416, and RQ-004845---have a different printed normal-closure
polynomial in the reconstruction.  Both records are retained.  Literal
polynomial equality is not used to infer field identity or packet
transport, and none of these discrepancies is used as a case-level
identity.  This is why mechanism eligibility remains separate from
higher-order packet proof.

\section{Reproducibility and relation to earlier work}

Every material assertion has a versioned JSON artifact, input hashes,
and replay code under PARI/GP 2.15.4.  From the project root, the local
revision gates are
\begin{center}\small
\texttt{python3 scripts/audit\_census\_paper.py}\quad and\quad
\texttt{python3 scripts/audit\_census\_referee\_revision.py}.
\end{center}
The quadratic theorem gate is
\texttt{python3 proof/audit\_q\_euler\_deleted\_prime\_cover.py}.
The primary finite evidence is the support-first reconciliation, the
row-level quadratic corpus, the deleted-prime cover certificate, and
the H-taxonomy matrix.  The deterministic version-1.0 paper-and-replay
package is archived at \url{https://doi.org/10.5281/zenodo.21730707}.
This version-1.1 correction candidate is local and is not represented
by that immutable version DOI.  The companion paper proves the uniform
quadratic theorem and selected higher-order identities; this census
does not reprove or enlarge them.

Table~\ref{tab:prior} states the exact comparison boundary.  No row
count is assigned to a prior theorem without a row-level map of its
extension, $S$-set, distinguished place, normalization, and conclusion.

\begin{table}[ht]
\centering\small
\caption{Prior-theorem boundary used in this census.}\label{tab:prior}
\begin{tabular}{>{\raggedright\arraybackslash}p{2.5cm}
                >{\raggedright\arraybackslash}p{3.5cm}
                >{\raggedright\arraybackslash}p{6.6cm}}\toprule
source & theorem/object & use and remaining boundary\\\midrule
Tate \cite{Tate1984} & Thm.~IV.5.4, rank one & Algebraicity input for
the quadratic formula after exact conductor, Euler, class-number,
unit-index, and orientation checks.\\
Dummit--Sands--Tangedal \cite{DST2003} & Thms.~2 and 4 & Refined
multiquadratic and full biquadratic antecedents; no census-row coverage
count is claimed without the still-open row-level hypothesis map.\\
Roblot \cite{Roblot2013} & Thms.~5.5, 6.1, 7.1 & Quadratic, cyclic
quartic, and conditional cyclic-sextic weak results.
Table~\ref{tab:higher} counts only exact Theorem~7.1 weak coverage,
never Artin-labelled packet identification.\\
Cohen--Roblot \cite{CohenRoblot2000} & Hilbert class-field algorithm &
The objects differ from one-place ray moduli.  Four exact subfield
controls test field containment only, not unit or packet overlap.\\\bottomrule
\end{tabular}
\end{table}

Arakawa's relative-index formula \cite{Arakawa1985} and the explicit
computational work of Dummit--Sands--Tangedal \cite{DST1997} are further
antecedents, but neither is assigned a census-row count here.

The prior-art comparison is deliberately narrow.  Cohen--Roblot study
Hilbert class fields, whereas the present promoted examples and census
rows are one-place ray-field objects with nontrivial finite modulus.
An exact four-control subfield audit is mixed: the selected ray normal
closures for bases $35$, $42$, and $186$ contain the associated Hilbert
biquadratic field, while the selected base-$51$ closure does not.
This is a PROVED statement about the four frozen normal closures only;
it identifies neither a Stark unit nor an Artin-labelled packet and
does not generalize containment.  This paper makes no priority assertion
from any of these comparisons.  The contribution claimed here is the
canonical finite classification and the exhaustive quadratic
packet-value-orbit corpus within its corrected boundary.

\appendix
\section{Record map and replay boundary}

\begin{sloppypar}
The trichotomy is in
\url{artifacts/census-paper-layer0-reconciliation-v1.json}; the
quadratic audit is in
\url{artifacts/census-q-packet-corpus-audit-v1.json}; the H matrix is
in \url{artifacts/census-h-taxonomy-v2.json}; the range-description
correction is in
\url{data/census-paper-preregistration-amendment-v18.json}; the preserved v1 record
used the legacy displayed-index proxy and is corrected by the v2
provenance fields.  The quadratic deleted-prime theorem and its
expanded-range counterexample are in
\url{artifacts/q-euler-deleted-prime-cover-theorem-v1.json}.
Engine-B transport
scope is in \url{artifacts/engine-b-transport-manifest-v5.json}; its
successor ledger is in \url{artifacts/engine-b-transport-ledger-v5.json}.
The latter records exactly twenty-two completed member transports.
The corrected global direct coprime-deletion geometry audit is recorded
in \url{artifacts/engine-b-global-coprime-geometry-audit-v1.json}; the
superseded B5-086/B5-021/B5-033 screens remain preserved with their
integral-basis correction record.
The exact four-control Hilbert/ray object audit is
in \url{artifacts/b5079-hilbert-ray-containment-v1.json} and
\url{artifacts/hilbert-ray-containment-tranche-v1.json}; it is field
containment only.  The versioned Cycle-127 correction record freezes
the revised source hashes, audit results, and replay resource use.  The
release record, including the root-file inventory and public checksums,
is associated with DOI
\url{https://doi.org/10.5281/zenodo.21730707}.
\end{sloppypar}

\begin{thebibliography}{20}
\raggedright
\footnotesize
\bibitem{ZhaoResults} H.~Zhao, \emph{Effective Archimedean Stark
Theorems over Real Quadratic Fields: Quadratic Support, Shintani
Transfer, and CM Descent}, version 1.5, 2026,
\url{https://doi.org/10.5281/zenodo.21713178}.
\bibitem{Tate1984} J.~Tate, \emph{Les conjectures de Stark sur les
fonctions $L$ d'Artin en $s=0$}, Progress in Mathematics 47,
Birkh\"auser, 1984.
\bibitem{Arakawa1985} T.~Arakawa, \emph{On the Stark--Shintani
conjecture and certain relative class numbers}, Adv. Stud. Pure Math.
7 (1985), 1--6.
\bibitem{DST1997} D.~S.~Dummit, J.~W.~Sands, and B.~A.~Tangedal,
\emph{Computing Stark units for totally real cubic fields}, Math.
Comp. 66 (1997), 1239--1267.
\bibitem{DST2003} D.~S.~Dummit, J.~W.~Sands, and B.~A.~Tangedal,
\emph{Stark's conjecture in multiquadratic extensions, revisited}, J.
Th\'eorie des Nombres de Bordeaux 15 (2003), 83--97.
\bibitem{CohenRoblot2000} H.~Cohen and X.-F.~Roblot, \emph{Computing
the Hilbert class field of real quadratic fields}, Math. Comp. 69
(2000), 1229--1244.
\bibitem{Roblot2013} X.-F.~Roblot, \emph{Index formulae for Stark
units and their solutions}, Pacific J. Math. 266 (2013), 391--422.
\end{thebibliography}
\end{document}
